\section{Findings, Limitations and Next Decisions}
The project connects instructor configuration to course-grounded responses and
proactive support through explicit evidence, state and execution contracts.
The comparison results identify three engineering findings. First, the
evidence-to-generator interface can matter more than increasing retrieval
complexity: typed targets improved a development comparison, while additional
ranking did not. Second, guarded planning and model allocation should be judged
against their stated objective: utility improved in the planner comparison
without improving every action label. Third, operational correctness and
instructional usefulness diverge: delivered, persistent check-ins can remain
generic. The objective-scoped correction improved completion correctness but
did not increase mean final simulator mastery relative to the old logic.

Flexible instruction, visual
retrieval and BKT/value timing have different development statuses. In
particular, the default planner still consumes delivery/event-based proxies
rather than a calibrated model of the committed learner evidence. The
objective-scoped completion correction addresses a real defect but retains
fixed thresholds and cumulative contradiction handling.

Validity is limited by synthetic courses, mechanical questions, simulator
assumptions and shared AI involvement in construction and review. Some
historical trials retain aggregates but lack the original per-case artifacts;
the appendix records those gaps rather than reconstructing evidence. Local
tests and virtual time do not establish production-scale uptime. The evaluations did not include an approved study with the actual instructor
or real students; those outcomes remain untested.

The next design decision is to compare a planner adapter based on committed
concept evidence against the present proxy, while retaining a simple control.
That comparison should then be integrated with source-supported,
profile-sensitive question generation and independently reviewed assessment.
An approved course pilot would test usefulness and student experience after
those gates. 